\chapter{Tinea corporis}
\index{Tinea corporis}

\section*{Definition and Overview}
Tinea corporis, colloquially known as ringworm, is a superficial cutaneous fungal infection of the glabrous (hairless) skin, excluding the scalp, beard, face, hands, feet, and groin. Despite its common name, the condition is not caused by any parasitic worm, but rather by microscopic fungi known as dermatophytes. These organisms have a unique affinity for keratinised tissues, such as the stratum corneum of the epidermis, hair, and nails. Tinea corporis is characterized by its classic annular (ring-like) lesions, which present with an active, erythematous, scaling border and a central clearing. This characteristic shape occurs as the fungus spreads centrifugally outward while the inflammatory response in the center subsides.

The significance of tinea corporis in clinical practice and public health is substantial. It is one of the most common dermatological complaints worldwide, affecting individuals of all ages, genders, and socioeconomic backgrounds. While generally non-life-threatening, tinea corporis can cause significant physical discomfort, including pruritus (itching) and localized pain, and can lead to secondary bacterial infections if the skin barrier is compromised by scratching. Furthermore, the cosmetic appearance of the lesions can cause significant psychological distress, social stigmatization, and anxiety for patients. Because dermatophytes are highly contagious, understanding the transmission dynamics, diagnostic modalities, and management strategies is critical for preventing outbreaks, particularly in closed environments like households, schools, gyms, and long-term care facilities. As global travel and participation in contact sports rise, the clinical burden of tinea corporis remains a key focus for dermatologists, primary care physicians, and public health practitioners alike.

\section*{Epidemiology}
The epidemiology of tinea corporis is marked by global prevalence, though the specific causative agents vary significantly by geographic location, climate, and local lifestyle factors. Tinea corporis occurs worldwide, but it is particularly prevalent in tropical and subtropical regions where high temperatures and humidity provide an optimal environment for fungal growth and proliferation. In these climates, the prevalence of superficial fungal infections can reach substantial proportions of the population.

In terms of demographics, tinea corporis can affect individuals at any stage of life, from neonates to the elderly. However, certain populations exhibit higher rates of infection. Children, particularly prepubertal children, are frequently affected due to close physical contact in school and childcare settings, as well as a developing immune response. Among adults, the infection is often associated with specific occupational or recreational activities. For instance, athletes engaged in contact sports---most notably wrestling, where the variant is termed ``tinea corporis gladiatorum''---show high rates of transmission due to skin-to-skin contact and shared equipment. Additionally, animal handlers, farmers, and pet owners are at a higher risk of contracting zoophilic species of dermatophytes.

No significant gender predilection exists for tinea corporis in the general population, though specific clinical settings, such as contact sports or occupational exposures, may show a higher prevalence in one gender over another. Socioeconomic factors also play a critical role; overcrowded living conditions, poor sanitation, and limited access to clean water or healthcare facilities significantly increase the risk of transmission and prolonged infection. Over the past few decades, the widespread use of over-the-counter topical corticosteroids has altered the clinical and epidemiological landscape, leading to an increase in atypical presentations (tinea incognito) and drug-resistant fungal strains, representing an emerging challenge in global dermatology.

\section*{Aetiology and Pathophysiology}
The development of tinea corporis involves a complex interplay between the infecting fungal pathogen, environmental conditions, and the host's cutaneous defenses.

\begin{itemize}
    \item \textbf{Causative Organisms:} Tinea corporis is caused by dermatophytes, which are filamentous fungi belonging to three primary genera: \textit{Trichophyton}, \textit{Microsporum}, and \textit{Epidermophyton}. The most common species isolated globally is \textit{Trichophyton rubrum}, an anthropophilic (human-associated) fungus that tends to cause chronic, recalcitrant infections. Other significant species include \textit{Trichophyton mentagrophytes}, \textit{Trichophyton interdigitale}, and \textit{Microsporum canis}. The latter is a zoophilic (animal-associated) fungus frequently transmitted to humans from infected domestic pets, such as cats and dogs, and typically induces a more intense inflammatory response.
    \item \textbf{Keratinase Production:} The hallmark of dermatophyte pathogenicity is the production of specialized extracellular enzymes known as keratinases. Since dermatophytes are unable to penetrate viable, living tissue in immunocompetent hosts, they colonize the non-living stratum corneum of the skin. Keratinases and other proteolytic enzymes degrade keratin into smaller peptides, which the fungus utilizes as nutrients, facilitating its invasion and spread across the skin surface.
    \item \textbf{Transmission Pathways:} Infection is acquired through contact with infectious arthroconidia (fungal spores) or hyphae. This occurs via three primary modes: anthropophilic transmission (direct skin-to-skin contact with an infected person or indirect contact with fomites like towels, clothing, or bed linens), zoophilic transmission (direct contact with infected animals), or geophilic transmission (direct contact with soil containing dermatophytes, such as \textit{Microsporum gypseum}).
    \item \textbf{Host Defense Mechanisms:} The human body employs several innate defenses against fungal invasion. These include the desquamation (shedding) of the stratum corneum, which physically removes colonizing fungi; the presence of sebum, which contains fungistatic fatty acids; and the activation of the innate immune system. Pattern recognition receptors on epidermal cells detect fungal cell wall components, triggering the release of pro-inflammatory cytokines and chemokines that recruit neutrophils and macrophages to the site.
    \item \textbf{Cell-Mediated Immunity:} A robust cell-mediated immune response, primarily driven by T-helper 1 (Th1) cells and interferon-gamma, is critical for clearing the infection. Individuals with impaired cell-mediated immunity, such as those with HIV/AIDS, patients undergoing immunosuppressive chemotherapy, or individuals taking systemic corticosteroids, are prone to more severe, widespread, or deep-seated dermatophyte infections, such as Majocchi's granuloma.
    \item \textbf{Risk Factors:} Several factors compromise the skin barrier or local immunity, predisposing individuals to tinea corporis. These include excessive sweating (hyperhidrosis), prolonged moisture retention in skin folds, minor skin trauma, occlusive clothing, diabetes mellitus, obesity, and pre-existing dermatological conditions like atopic dermatitis.
\end{itemize}

\section*{Clinical Features}
The clinical presentation of tinea corporis is characteristic, but can vary depending on the causative organism, the host's immune status, and any prior topical treatments.

\begin{itemize}
    \item \textbf{Annular Lesions:} The classic presentation of tinea corporis is one or more circular or oval erythematous plaques. As the lesion expands centrifugally, the center typically clears, leaving a ring-like appearance. The active, advancing border is raised, erythematous, and often scaling. In some cases, vesicles or pustules may be present along this active edge, indicating a more intense inflammatory response.
    \item \textbf{Pruritus (Itching):} Intense itching is the most common and troublesome symptom reported by patients. The severity of pruritus can vary, but it frequently leads to excoriation (skin scratching), which increases the risk of secondary bacterial infections and lichenification (thickening of the skin).
    \item \textbf{Plaque Distribution and Size:} Lesions can occur on any part of the glabrous skin, including the trunk, limbs, and neck. They can range in size from a few millimeters to several centimeters in diameter. Multiple lesions may coalesce, forming large, polycyclic, or geographic patterns across the skin surface.
    \item \textbf{Tinea Incognito:} If a patient has treated the lesion with topical corticosteroids, the typical inflammatory features and scaling border may be suppressed, while the fungus continues to proliferate. This results in an atypical, poorly defined, non-scaling lesion that is difficult to diagnose clinically, a phenomenon known as tinea incognito.
    \item \textbf{Majocchi's Granuloma:} This is a deep-seated dermatophyte infection that occurs when the fungus penetrates the hair follicle into the dermis, often secondary to shaving or the use of topical steroids. It presents as erythematous, indurated nodules or pustules, typically on the lower legs of females, and requires systemic rather than topical antifungal therapy.
    \item \textbf{Inflammatory vs. Non-inflammatory Variants:} Anthropophilic species like \textit{Trichophyton rubrum} often produce mild, minimally inflammatory scaling plaques that can persist for months. In contrast, zoophilic species like \textit{Microsporum canis} or \textit{Trichophyton verrucosum} typically trigger an acute, highly inflammatory reaction with pustules, vesiculation, and significant discomfort.
\end{itemize}

\section*{Differential Diagnosis}
Because many dermatological conditions present with annular, scaling, or erythematous lesions, tinea corporis must be carefully differentiated from other disorders to avoid inappropriate treatment.

\begin{itemize}
    \item \textbf{Nummular Eczema:} Also known as discoid eczema, this condition presents with coin-shaped, highly pruritic, erythematous plaques. Unlike tinea corporis, nummular eczema typically lacks central clearing, is often oozing or crusting in its acute phase, and does not show fungal elements on microscopic examination.
    \item \textbf{Plaque Psoriasis:} Psoriatic lesions are characterized by well-demarcated, erythematous plaques covered with thick, silvery scales. They commonly affect the extensor surfaces (elbows, knees) and the scalp. Removal of the scale may reveal pinpoint bleeding (Auspitz's sign), and psoriasis lacks the active, scaling advancing edge with central clearing seen in tinea corporis.
    \item \textbf{Pityriasis Rosea:} This self-limiting inflammatory skin disease classically begins with a single, larger ``herald patch'' on the trunk, followed days to weeks later by a widespread eruption of smaller, oval, scaling plaques distributed along the Langer's lines of the back (often described as a ``Christmas tree'' pattern). The scale is typically collarette (attached at the periphery and loose in the center), distinguishing it from ringworm.
    \item \textbf{Contact Dermatitis:} Both allergic and irritant contact dermatitis can present with itchy, red, scaling plaques. However, these lesions are typically confined to the area of contact with the offending substance (e.g., jewelry, cosmetics, or chemicals) and lack the symmetrical, expanding annular configuration of tinea corporis.
    \item \textbf{Granuloma Annulare:} This benign inflammatory skin condition presents with annular, smooth, non-scaling dermal papules or plaques arranged in a ring. Because there is no epidermal involvement, there is no scaling, which is a key clinical differentiator from tinea corporis.
    \item \textbf{Subacute Cutaneous Lupus Erythematosus (SCLE):} SCLE can present with annular or polycyclic, scaling, erythematous lesions, primarily in sun-exposed areas. It is associated with systemic autoimmune markers (such as anti-Ro/SSA antibodies) and histopathological features of interface dermatitis, rather than fungal hyphae.
    \item \textbf{Erythema Annulare Centrifugum (EAC):} EAC is a reactive dermatosis characterized by annular lesions that spread outwards with a characteristic ``trailing scale'' inside the advancing erythematous border, unlike the leading scale seen on the outermost edge in tinea corporis.
\end{itemize}

\section*{When to Seek Medical Care}
While tinea corporis is generally a superficial and self-limiting infection that responds well to over-the-counter or prescription topical treatments, certain circumstances warrant evaluation by a healthcare provider. Patients should consult a primary care physician or a dermatologist if:
\begin{itemize}
    \item The lesions do not show improvement after two weeks of consistent treatment with over-the-counter antifungal creams.
    \item The rash covers a large surface area of the body or appears to be spreading rapidly despite treatment.
    \item There are signs of a secondary bacterial infection, such as increased pain, warmth, swelling, red streaks spreading from the lesion, or the presence of purulent discharge (pus).
    \item The patient has a compromised immune system (e.g., due to HIV/AIDS, cancer treatments, or immunosuppressive medications) or has diabetes, as they are at a higher risk of systemic complications.
    \item The lesions are located on the scalp or face, or involve the nails, as these areas typically require oral prescription antifungal medications.
\end{itemize}

While tinea corporis itself does not cause life-threatening emergencies, systemic infections or severe allergic reactions (anaphylaxis) to antifungal medications require immediate emergency medical attention. If a patient experiences severe symptoms such as difficulty breathing, swelling of the face, lips, or throat, or a sudden, widespread rash accompanied by high fever, they must seek emergency services immediately. In many countries, call local emergency services; in the United States or Canada, call 911; in the United Kingdom, call 999; or contact the local emergency service number.

\section*{Diagnosis and Investigations}
Confirming the diagnosis of tinea corporis is essential to ensure appropriate treatment and avoid the misuse of topical corticosteroids. While a clinical diagnosis is often made based on the characteristic appearance of the lesions, laboratory investigations are invaluable, particularly in atypical, severe, or treatment-resistant cases.

\begin{itemize}
    \item \textbf{Physical Examination:} A thorough dermatological examination is the first step. The clinician assesses the distribution, shape, border, and scaling pattern of the lesions. Using a handheld magnifying lens or a dermatoscope can help visualize the fine scale at the border and rule out other annular dermatoses.
    \item \textbf{Potassium Hydroxide (KOH) Preparation:} This is the most rapid, cost-effective, and widely used diagnostic test. The clinician gently scrapes the active scaling edge of the lesion with a scalpel blade or the edge of a glass slide. The scraped scales are placed on a slide, treated with a 10\% to 20\% KOH solution (sometimes with dimethyl sulfoxide or a fungal stain like Chlorazol Black E), and gently heated. The KOH dissolves the keratin but leaves the fungal cell walls intact. Under light microscopy, the presence of branching, septate hyphae confirms a dermatophyte infection.
    \item \textbf{Fungal Culture:} Although it takes several weeks (usually 2 to 4 weeks) to obtain results, a fungal culture remains a gold standard for identifying the specific species of dermatophyte. Scrapings are inoculated onto selective media, such as Sabouraud's dextrose agar containing cycloheximide and chloramphenicol to suppress contaminating molds and bacteria. Identification of the species is helpful in tracing the source of infection (e.g., zoophilic source from a pet) and directing epidemiological investigations.
    \item \textbf{Wood's Lamp Examination:} A Wood's lamp emits ultraviolet light at a wavelength of approximately 365 nm. While most causative agents of tinea corporis (such as \textit{Trichophyton rubrum}) do not fluoresce, certain species of \textit{Microsporum} (like \textit{Microsporum canis}) may exhibit a bright blue-green fluorescence. This tool is particularly useful for screening pets or detecting concurrent tinea capitis (scalp ringworm).
    \item \textbf{Polymerase Chain Reaction (PCR) Testing:} Molecular diagnostics, such as PCR assays, are increasingly used in advanced clinical settings. PCR can detect dermatophyte DNA directly from skin scrapings within hours. It offers high sensitivity and specificity and can identify the exact species, though its higher cost and limited availability restrict its routine use.
    \item \textbf{Skin Biopsy:} In rare, atypical presentations, or when the diagnosis remains elusive after negative KOH preps and cultures, a skin biopsy may be performed. Histopathological examination with special stains, such as Periodic acid-Schiff (PAS) or Grocott's methenamine silver (GMS), reveals fungal hyphae in the stratum corneum, along with epidermal spongiosis and a superficial perivascular inflammatory infiltrate.
\end{itemize}

\section*{Management}
The management of tinea corporis involves a combination of topical or systemic antifungal pharmacotherapy and supportive lifestyle modifications to eradicate the pathogen, alleviate symptoms, and prevent recurrence.

\begin{itemize}
    \item \textbf{Topical Antifungal Therapy:} For localized, uncomplicated lesions of tinea corporis, topical antifungals are the first-line treatment. These are typically applied to the affected area and a 2 cm margin of normal-appearing skin once or twice daily for 2 to 4 weeks. Commonly used classes include:
    \begin{itemize}
        \item \textit{Allylamines:} Terbinafine 1\% cream and naftifine are highly effective. Allylamines are fungicidal (they kill the fungus) by inhibiting the enzyme squalene epoxidase, leading to a accumulation of toxic squalene and depletion of ergosterol in the fungal cell membrane. They often require shorter treatment courses (e.g., 1 to 2 weeks).
        \item \textit{Imidazoles:} Clotrimazole 1\%, miconazole 2\%, ketoconazole 2\%, and econazole are fungistatic (they inhibit fungal growth) by blocking the enzyme 14-alpha-demethylase. They require a longer application period (usually 3 to 4 weeks) and compliance is critical to prevent relapse.
        \item \textit{Other Agents:} Ciclopirox olamine 1\% and amorolfine are alternative topical options with broad-spectrum antifungal and anti-inflammatory properties.
    \end{itemize}
    \item \textbf{Systemic Antifungal Therapy:} Oral antifungal medications are indicated for extensive or widespread infections, Majocchi's granuloma, cases involving immunocompromised patients, or when topical treatment has failed. Before initiating systemic therapy, liver function tests may be recommended, particularly for prolonged courses. Key oral agents include:
    \begin{itemize}
        \item \textit{Terbinafine:} Standard adult dose is 250 mg once daily for 1 to 2 weeks. It is highly effective and achieves high concentrations in the stratum corneum and hair follicles.
        \item \textit{Itraconazole:} Given at 100 to 200 mg daily for 1 to 2 weeks. It is a broad-spectrum triazole that requires gastric acidity for optimal absorption.
        \item \textit{Fluconazole:} Typically dosed at 150 mg to 200 mg once weekly for 2 to 4 weeks.
        \item \textit{Griseofulvin:} Historically common, but now largely superseded by newer agents due to longer treatment courses (4 to 6 weeks) and higher rates of side effects and drug interactions.
    \end{itemize}
    \item \textbf{Symptomatic Relief:} Pruritus can be managed with mild, non-sedating oral antihistamines. The use of topical combination products containing an antifungal and a potent corticosteroid (e.g., clotrimazole-betamethasone) is generally discouraged. While they may provide rapid temporary relief from itching, the steroid component can suppress local immunity, promote fungal replication, mask symptoms, and increase the risk of treatment failure or conversion to tinea incognito.
    \item \textbf{Environmental Sanitisation:} To prevent reinfection, patients should wash all clothing, towels, and bed linens in hot water (at least 60 degrees Celsius) and dry them on high heat. Personal items must not be shared. If a zoophilic source is suspected, domestic pets should be evaluated and treated by a veterinarian.
\end{itemize}

\section*{Complications}
Although tinea corporis is generally a benign, superficial condition, complications can arise, especially if the infection is misdiagnosed, untreated, or occurs in vulnerable patient populations.

\begin{itemize}
    \item \textbf{Secondary Bacterial Infection:} Pruritus leads to intense scratching, which breaks the skin barrier and introduces opportunistic bacteria, most commonly \textit{Staphylococcus aureus} or \textit{Streptococcus pyogenes}. This can result in impetiginisation, cellulitis, or lymphangitis, requiring antibiotic therapy.
    \item \textbf{Tinea Incognito:} The inappropriate application of topical corticosteroids alters the appearance of the rash, making clinical diagnosis difficult and allowing the fungus to spread unchecked into deeper cutaneous structures.
    \item \textbf{Majocchi's Granuloma:} Fungal elements can bypass the epidermis and invade the dermis, causing a chronic, granulomatous folliculitis. This complication requires systemic antifungal therapy and can lead to scarring and localized alopecia (hair loss).
    \item \textbf{Hyperpigmentation and Hypopigmentation:} Following the resolution of the inflammatory phase of the infection, patients may experience post-inflammatory pigmentary changes. The skin may appear darker (hyperpigmentation) or lighter (hypopigmentation) than the surrounding tissue. While these changes are usually temporary, they can take several months to resolve, causing cosmetic concern.
    \item \textbf{Erythema Nodosum and Dermatophytid (``Id'') Reactions:} In rare instances, an intense inflammatory response to the dermatophyte can trigger systemic immunologic reactions. An ``id'' reaction is a secondary, sterile, pruritic skin eruption occurring at a site distant from the primary infection. Erythema nodosum, characterized by painful, red nodules on the shins, has also been reported as a rare immunological complication.
    \item \textbf{Psychosocial Impact:} Chronic or widespread tinea corporis can cause significant anxiety, low self-esteem, and social withdrawal due to the visible nature of the lesions and the incorrect stigma that the infection is due to poor personal hygiene.
\end{itemize}

\section*{Prognosis and Outcomes}
The prognosis for individuals with tinea corporis is excellent. In the vast majority of cases, the infection is superficial, localized, and responds promptly to appropriate topical antifungal therapy. Complete resolution of the lesions, scaling, and itching is typically achieved within two to four weeks of initiating treatment.

For patients requiring systemic antifungal therapy, the cure rates remain high. However, the time course to complete clinical clearance can be longer than the duration of active treatment, as the skin requires time to shed the damaged cells and regenerate a healthy stratum corneum. Post-inflammatory hyperpigmentation or hypopigmentation can persist for several weeks or months after the fungus has been eradicated, but these changes are benign and gradually fade without permanent scarring, unless a deep secondary bacterial infection or Majocchi's granuloma occurred.

Factors that can negatively affect the prognosis or lead to chronic, recurrent infections include:
\begin{itemize}
    \item Non-compliance with the prescribed treatment regimen (e.g., stopping the antifungal cream as soon as the itching stops, rather than completing the full recommended course).
    \item Ongoing exposure to a persistent source of infection, such as an untreated household member, a contaminated environment, or an infected pet.
    \item Underlying medical conditions that compromise the host's immune system, such as poorly controlled diabetes mellitus, peripheral vascular disease, or systemic immunosuppressive therapy.
    \item The presence of drug-resistant dermatophyte strains, which have emerged in recent years, particularly in parts of South Asia, due to the widespread availability and misuse of fixed-dose combination creams containing high-potency steroids, antibacterials, and antifungals. These cases require prolonged, culture-guided oral therapy and close dermatological follow-up.
\end{itemize}

\section*{Prevention and Self-Care}
Preventing the acquisition and spread of tinea corporis involves practical self-care measures, strict hygiene protocols, and environmental modifications aimed at reducing exposure to dermatophyte spores.

\begin{itemize}
    \item \textbf{Maintain Clean, Dry Skin:} Dermatophytes thrive in warm, moist environments. After bathing, swimming, or exercising, dry the skin thoroughly, paying particular attention to skin folds, the groin, and the areas between the toes.
    \item \textbf{Practice Good Hand Hygiene:} Wash hands thoroughly with soap and water after touching pets, working in the garden, or coming into contact with anyone suspected of having a fungal infection.
    \item \textbf{Avoid Sharing Personal Items:} Do not share clothing, towels, hairbrushes, bed linens, sports gear, or other personal grooming items, as these can harbor infectious fungal spores for extended periods.
    \item \textbf{Wear Loose, Breathable Clothing:} Choose clothing made of natural fibers, such as cotton or linen, which allow moisture to evaporate. Avoid tight-fitting or synthetic garments that trap sweat and heat against the skin.
    \item \textbf{Protect Skin in Public Facilities:} Wear shower shoes, sandals, or flip-flops in communal showers, locker rooms, public pools, and gym changing areas to avoid direct contact with contaminated floors.
    \item \textbf{Disinfect Sports Equipment:} Clean and disinfect exercise mats, helmets, and other shared sports gear before and after use. Athletes in contact sports should shower immediately after training or competition.
    \item \textbf{Manage Pet Health:} Regularly inspect domestic pets for signs of ringworm, such as patches of hair loss, scaling, or redness. If a pet shows symptoms, seek prompt veterinary evaluation and treatment. Wash pet bedding regularly.
    \item \textbf{Complete the Treatment Course:} Always use prescribed or over-the-counter antifungal medications for the full duration recommended by the clinician or product packaging, even if the symptoms disappear early. This ensures complete eradication of the fungus and minimizes the risk of recurrence.
\end{itemize}

\section*{Sources and Further Reading}
For authoritative information, diagnostic guidelines, and updates on tinea corporis and other superficial fungal infections, refer to the following resources:
\begin{itemize}
    \item \textbf{World Health Organization (WHO):} Guidelines on the management of common skin diseases, including superficial fungal infections. URL: https://www.who.int/
    \item \textbf{Centers for Disease Control and Prevention (CDC):} Detailed patient education and clinical resources on dermatophyte infections (ringworm). URL: https://www.cdc.gov/fungal/diseases/ringworm/
    \item \textbf{National Center for Biotechnology Information (NCBI) - StatPearls:} Comprehensive medical review of tinea corporis, its diagnosis, and management. URL: https://www.ncbi.nlm.nih.gov/books/NBK541075/
    \item \textbf{DermNet NZ:} An authoritative, peer-reviewed clinical resource specializing in dermatology and skin conditions. URL: https://dermnetnz.org/topics/tinea-corporis/
    \item \textbf{American Academy of Dermatology (AAD):} Public health advisories, skin care tips, and dermatological guidelines for patients and practitioners. URL: https://www.aad.org/public/diseases/a-z/ringworm/
    \item \textbf{British Association of Dermatologists (BAD):} Patient information leaflets and clinical guidelines on tinea infections. URL: https://www.bad.org.uk/
\end{itemize}